\section{Introduction}\label{sec:intro}

Federated learning (FL) \cite{pmlr-v54-mcmahan17a} enables collaborative model training across distributed clients without exchanging raw data, making it a promising paradigm for privacy-sensitive visual learning. Despite this advantage, practical FL deployments must balance three competing factors: statistical heterogeneity across clients \cite{tan2022fedproto,10660707,nguyen2026towards}, limited computational resources on edge devices \cite{MAL-083,11180918}, and communication efficiency in large-scale distributed systems \cite{nguyen2022hcfl,nguyen2025federated}. Conventional FL methods typically require each client to train the full model locally, which can be prohibitively expensive for resource-constrained devices \cite{9760729,9084352,dai2023tacklingFedProx}.

Split federated learning (SFL) addresses this limitation by partitioning model execution across clients and servers \cite{thapa2022splitfed}. By assigning only shallow network layers to clients, SFL reduces client-side computation and memory requirements \cite{GUPTA20181,vepakomma2019split}. Hierarchical SFL further extends this idea to cloud--edge--client systems, where edge servers provide intermediate computation and aggregation. However, existing SFL and hierarchical SFL (HSFL) frameworks usually rely on transmitting intermediate activations during the forward pass and propagating gradients during backpropagation \cite{11091500,zhang2025amperecommunicationefficienthighaccuracysplit,qiu2023evaluating,wang2025towards}. This sample-level communication creates substantial bandwidth consumption, introduces synchronization dependencies across tiers, and couples the optimization of client, edge, and cloud components \cite{10980018,quan2023hiersfl}. These challenges become more pronounced as models, client populations, and data heterogeneity increase.

In this work, we revisit the communication interface of hierarchical split learning. Instead of exchanging sample-level activations and gradients, we ask whether compact class-level feature statistics can provide sufficient information for training upper model components. Our key observation is that first- and second-order statistics of client-side feature embeddings can preserve useful class-level semantic structure while requiring much less communication than activation-level exchange. This motivates a statistical feature-prototyping mechanism, in which upper tiers synthesize representative feature distributions from compact summaries rather than relying on direct cross-tier backpropagation.

Based on this principle, we propose \textbf{Hierarchical Split-Federated Prototyping (H-SFP)}, a communication-efficient framework for cloud--edge--client learning. In H-SFP, clients execute lightweight feature extractors and transmit only class-wise feature statistics to nearby edge servers. Edge servers aggregate these statistics, synthesize representative feature samples, and train intermediate model components. The cloud tier further aggregates distribution-level summaries to train task-specific layers. By replacing activation and gradient exchange with statistical feature prototyping, H-SFP preserves the computational advantage of split-model execution while decoupling optimization across hierarchical tiers and reducing communication overhead.

The main finding of this work is that compact statistical summaries can serve as an effective communication interface for HSFL. Although H-SFP does not transmit sample-level activations or gradients across tiers, it maintains strong performance under heterogeneous data distributions while substantially reducing communication cost. Our contributions are summarized as follows:
\begin{itemize}
\item We propose H-SFP, a communication-efficient hierarchical split-federated framework that replaces sample-level activation and gradient exchange with compact class-wise feature statistics.

\item We introduce a statistical feature-prototyping mechanism that enables the edge and cloud tiers to synthesize representative feature distributions and to train upper-model components without cross-tier backpropagation.

\item We provide theoretical analysis of prototype-induced distribution shift and the stability of statistical feature representations in hierarchical setups.

\item We conduct experiments on image classification and medical image segmentation benchmarks, demonstrating robust performance under heterogeneous data distributions with substantially reduced communication overhead.
\end{itemize}
